\begin{abstract}
Attention is widely used as a proxy for visual grounding in large vision--language models (LVLMs): detectors score how object tokens attend to the image, while mitigation methods amplify or redistribute visual attention to reduce hallucinations. We argue that this premise is incomplete: visual attention allocation is not object-level evidence verification. Attention can be plausibly but incorrectly grounded, routing to regions semantically associated with the target object without verifying the target itself. We propose a verification-centered diagnostic framework for testing this proxy in both detection and mitigation. On the detection side, we show that strong attention-based AUROC can be driven by \emph{generation-position confounding}. On LLaVA-1.5-7B over MSCOCO ($16{,}426$ object mentions), position alone reaches $0.830$ AUROC; under within-bin, matched-pair, same-object, and residual evaluation, attention-mass detectors such as PAS and SVAR retain only a small fraction of their above-chance signal, while representation and uncertainty baselines remain more robust. A purpose-built attention-shape stress test, SinkDetect, also collapses under position control. On the mitigation side, we reproduce attention-intervention components of PAI, ClearSight, and Visual Attention Sink on POPE and CHAIR. ClearSight mainly shifts POPE answers toward ``yes'' (TPR $+3.5$ points, FPR $+3.8$ points), while an attention-shift audit shows that PAI increases active-layer visual attention by $+0.038$ without changing any adversarial POPE decisions. Finally, a semantic-neighbor POPE audit shows that all attention-only interventions have substantially higher false-positive rates when related objects are visible than when the target is plainly absent. As a constructive check, we formalize Target-Discriminative Evidence Verification (TDEV) and instantiate it with OWLv2 region evidence: the verifier reaches $0.849$ within-bin and $0.851$ matched-pair AUROC on MSCOCO mention detection, and a two-stage POPE gate lowers macro related-present FPR to $0.069$ while retaining macro MCC $0.751$. These results suggest that attention can contain grounding information, but simple attention scores and attention amplification do not by themselves verify whether a visual claim is supported; target-discriminative region evidence is a more reliable direction.
\end{abstract}
